% Chapter 14: Connections and Applications
% Part IV: Implementation and Applications

\chapter{연결과 응용}
\label{ch:connections}

\begin{quote}
\textit{``좋은 이론은 고립되지 않는다. 기존 지식의 지형 속에서 자신의 위치를 밝히고, 새로운 탐험의 방향을 제시한다.''}
\end{quote}

\noindent
SCC 이론은 수학적으로 자기-완결적이지만, 그 가치는 다른 학문과의 연결을 통해 배가된다. 이 장에서는 게슈탈트 심리학, 예측적 처리, 정보 통합 이론(IIT), 베이지안 뇌, 능동 추론과의 관계를 탐구하고, 10개의 경험적 예측을 제시하며, 컴퓨터 비전과 인지 모델링에의 응용 가능성을 논의한다.


%% ============================================================
\section{SCC와 게슈탈트 심리학}
\label{sec:gestalt}

게슈탈트 심리학의 조직화 법칙(laws of organization)은 지각적 그룹화의 현상학적 기술로, 100년 이상의 역사를 가진다. SCC는 이 법칙들에 대한 \textbf{형식적 실현}을 제공한다---각 게슈탈트 원리가 SCC의 어떤 수학적 구성 요소에 대응하는지를 정밀하게 밝힌다.

\begin{table}[h]
\centering
\caption{게슈탈트 법칙과 SCC 대응표}
\label{tab:gestalt-mapping}
\begin{tabular}{lll}
\toprule
\textbf{게슈탈트 법칙} & \textbf{SCC 대응} & \textbf{형식적 근거} \\
\midrule
근접성 (Proximity) & 인접 행렬 $N(x,y)$ & 공간적 가중치 \\
유사성 (Similarity) & 응집 가중 전달 비용 & 지문 $\varphi$ 기반 비용 \\
폐합 (Closure) & $\mathrm{Cl}_t(u)$ & 자기-완성 연산자 \\
연속 (Good Continuation) & $\mathcal{E}_{\mathrm{bd}}$ 평활도 항 & $2\alpha u^T L u$ \\
공통 운명 (Common Fate) & 시간적 전달 $\mathbf{M}_{t \to s}$ & 구조적 계승 \\
Pr\"agnanz (간결성) & 에너지 최소화 & $\hat{u} = \arg\min \mathcal{E}$ \\
\bottomrule
\end{tabular}
\end{table}

\subsection{근접성 $\to$ 인접 구조}

게슈탈트의 근접성 법칙은 ``가까운 것들이 함께 묶인다''는 것이다. SCC에서 이는 인접 행렬 $N(x,y)$의 구조로 자연스럽게 포착된다. 4-연결 격자에서 $N(x,y) = 1$이면 $x$와 $y$는 이웃이며, $P_t u(x)$는 $x$의 이웃에서의 평균 응집을 나타낸다. 근접한 노드가 높은 $u$ 값을 공유할 때, 폐합 연산자가 이를 강화한다.

그러나 SCC는 단순 근접성을 넘어선다. 인접 행렬을 일반 그래프로 확장하면, 근접성은 공간적 거리가 아닌 \emph{관계적 거리}로 일반화된다. 이는 의미적 근접성(semantic proximity)이나 기능적 연결(functional connectivity)도 포착할 수 있게 한다.

\subsection{유사성 $\to$ 응집 가중 전달 비용}

게슈탈트의 유사성 법칙은 ``비슷한 것들이 함께 묶인다''는 것이다. SCC에서 유사성은 응집 지문 $\varphi(x) = (u(x), \mathrm{Cl}(u)(x), D(x; 1{-}u))$를 통해 포착된다. 지문이 유사한 노드 쌍은 낮은 전달 비용을 가지며, 이는 시간적 전달 $\mathbf{M}_{t \to s}$에서 높은 결합 질량으로 표현된다.

\subsection{폐합 $\to$ 자기-완성}

가장 직접적인 대응이다. 게슈탈트의 폐합 법칙---``불완전한 형태는 완성되려는 경향이 있다''---은 $\mathrm{Cl}_t$의 수학적 성질 그 자체이다. 시그모이드 폐합은 이웃의 높은 응집을 자신에게 전파하여, 부분적으로 형성된 패턴을 자기-강화한다. A3 축약 조건($a_{\mathrm{cl}} < 4$)은 이 자기-강화가 안정적으로 수렴함을 보장한다.

\subsection{연속 $\to$ 경계 평활도}

게슈탈트의 좋은 연속(good continuation) 법칙은 SCC에서 경계 에너지의 평활도 항 $2\alpha u^T L u$로 포착된다. 이 항은 응집 장의 급격한 변화를 억제하여, 부드러운 경계 전이를 선호한다. $\alpha / \beta$ 비율이 형성의 경계 두께를 제어한다.

\subsection{공통 운명 $\to$ 시간적 전달}

공통 운명 법칙---``같은 방향으로 움직이는 것들이 함께 묶인다''---은 시간적 전달 $\mathbf{M}_{t \to s}$의 동기에 직결된다. 동일 형성에 속하는 노드들은 유사한 지문을 가지므로, 시간 전달에서 일관되게 대응된다. T-Persist-1(e)의 핵심-대-핵심 집중은 이를 정량적으로 보장한다.

\paragraph{게슈탈트를 넘어서.} SCC는 게슈탈트 법칙을 단순히 재진술하는 것이 아니다. 세 가지 면에서 확장한다:
\begin{enumerate}
    \item \textbf{정량화}: 각 원리의 상대적 강도가 에너지 가중치 $\lambda_i$로 조절 가능
    \item \textbf{상호작용}: 원리들이 에너지 최소화를 통해 자동으로 경쟁/협력
    \item \textbf{예측}: 원리들의 상호작용에서 새로운 현상 예측 (예: Sep-before-Inside 순서, P5)
\end{enumerate}

% Figure placeholder
\begin{figure}[t]
\centering
\fbox{\parbox{0.85\textwidth}{\centering\vspace{5cm}
\textbf{[그림 14.1]} 게슈탈트 법칙과 SCC 대응 도표.\\
각 게슈탈트 원리(좌)가 SCC의 수학적 구성 요소(우)에 화살표로 연결.\\
다대다 대응: 근접성과 유사성은 모두 전달 비용에 기여.
\vspace{0.5cm}}}
\caption{게슈탈트 심리학과 SCC 이론의 형식적 대응. 각 원리에 대해 SCC는 정량적 실현과 상호작용 메커니즘을 제공한다.}
\label{fig:gestalt-mapping}
\end{figure}


%% ============================================================
\section{예측적 처리와의 비교}
\label{sec:predictive-processing}

예측적 처리(predictive processing, PP)는 뇌를 ``예측 기계''로 보는 프레임워크로, 감각 입력과 내부 모델의 불일치(예측 오류)를 최소화하는 것이 인지의 핵심이라고 주장한다.

\subsection{수렴점}

SCC와 PP는 여러 구조적 유사성을 공유한다:
\begin{itemize}
    \item \textbf{하향식 영향}: 폐합 연산자 $\mathrm{Cl}_t$는 전체적 응집 구조가 국소 값에 영향을 미치는 하향식(top-down) 메커니즘이다. 이는 PP의 예측이 감각 처리에 미치는 하향식 영향과 구조적으로 유사하다.
    \item \textbf{에너지 최소화}: SCC의 에너지 최소화는 PP의 예측 오류 최소화와 형식적으로 유사한 변분적 구조를 공유한다.
    \item \textbf{자기-참조성}: SCC의 에너지는 자기-참조적이다(폐합과 구별이 모두 $u$ 자체에 의존). PP에서 예측은 감각 데이터에 의해 갱신되고, 갱신된 예측이 다시 감각 처리에 영향을 미친다.
\end{itemize}

\subsection{발산점}

\begin{itemize}
    \item \textbf{존재론적 수준}: PP는 표상(representation)과 추론(inference)의 언어를 사용한다. SCC는 전-객체적(pre-objective) 수준에서 작동하며, 표상 개념을 전제하지 않는다.
    \item \textbf{시간적 구조}: PP의 예측-오류 역학은 초 단위의 빠른 시간 척도에서 작동한다. SCC의 시간적 전달은 형성의 구조적 계승을 다루며, 시간 척도가 특정되지 않는다.
    \item \textbf{계층}: PP는 본질적으로 계층적이다 (저수준 예측 → 고수준 예측). SCC는 단일 층위에서 작동하며, 계층적 확장은 미래 과제이다.
\end{itemize}


%% ============================================================
\section{IIT, 베이지안 뇌, 능동 추론과의 비교}
\label{sec:other-frameworks}

\subsection{정보 통합 이론 (IIT)}

IIT의 핵심 지표 $\Phi$는 시스템의 통합 정보를 측정한다. SCC의 진단 벡터 $\mathbf{d} = (\mathsf{Bind}, \mathsf{Sep}, \mathsf{Inside}, \mathsf{Persist})$와 구조적으로 비교할 수 있다:

\begin{itemize}
    \item $\Phi > 0$ $\leftrightarrow$ $\mathsf{Bind} > 0$: 둘 다 ``부분들이 전체로 통합되어 있는가?''를 묻는다
    \item IIT의 배제 공리(exclusion postulate) $\leftrightarrow$ SCC의 체적 제약 $\Sigma_m$: 유한한 자원을 경쟁적으로 배분
    \item 차이점: IIT는 의식의 존재/부재를 판별하려 하지만, SCC는 형성의 \emph{품질}을 연속적으로 측정한다
\end{itemize}

\subsection{베이지안 뇌}

베이지안 프레임워크에서 지각은 사전 분포(prior)와 우도(likelihood)의 곱을 최대화하는 사후 분포(posterior)이다. SCC와의 대응:

\begin{itemize}
    \item 사전 분포 $\leftrightarrow$ 경계 에너지 $\mathcal{E}_{\mathrm{bd}}$ (형성의 평활성과 이진성 선호)
    \item 우도 $\leftrightarrow$ 폐합/분리 에너지 (관찰된 응집 패턴과의 적합)
    \item 사후 분포 $\leftrightarrow$ 에너지 최소화자 $\hat{u}$
    \item 차이점: SCC에는 명시적 데이터(관찰)가 없다---에너지 자체가 자기-참조적이다
\end{itemize}

\subsection{능동 추론 (Active Inference)}

능동 추론에서 유기체는 자유 에너지(variational free energy)를 최소화한다. SCC의 에너지 $\mathcal{E}$와 변분 자유 에너지의 구조적 유사성:

\begin{itemize}
    \item 둘 다 ``정확성''(데이터 적합)과 ``복잡성''(모델 단순성) 사이의 트레이드오프를 포함
    \item SCC에서: 정확성 $\sim \mathcal{E}_{\mathrm{cl}} + \mathcal{E}_{\mathrm{sep}}$, 복잡성 $\sim \mathcal{E}_{\mathrm{bd}}$
    \item 차이점: 능동 추론은 행동(action)을 포함하지만, SCC는 순수 지각적이다
\end{itemize}

% Figure placeholder
\begin{figure}[t]
\centering
\fbox{\parbox{0.85\textwidth}{\centering\vspace{5cm}
\textbf{[그림 14.2]} 프레임워크 비교 매트릭스.\\
SCC, PP, IIT, 베이지안 뇌, 능동 추론의\\
5개 차원(자기-참조성, 연속성, 시간성, 계층성, 의식 여부)에서의 비교.
\vspace{0.5cm}}}
\caption{SCC와 네 가지 인지과학 프레임워크의 구조적 비교. SCC의 고유한 기여는 전-객체적 수준에서의 자기-참조적 에너지 최소화이다.}
\label{fig:framework-comparison}
\end{figure}


%% ============================================================
\section{열 가지 경험적 예측}
\label{sec:ten-predictions}

SCC 이론은 심리물리학, 신경영상, 발달심리학에서 검증 가능한 10개의 경험적 예측을 제시한다. 각 예측에 대해 검증 가능성 등급(A: 현재 기술로 직접 검증 가능, B: 가능하나 복잡한 설계 필요, C: 새로운 실험 패러다임 필요)을 부여한다.

\subsection{심리물리학 예측}

\paragraph{E1. 폐합 선행 효과 (등급 A).}
SCC는 폐합이 구별보다 에너지적으로 먼저 달성됨을 예측한다 (P5의 인지적 대응). 불완전한 형태의 지각적 완성이 내부/외부 판별보다 시간적으로 선행해야 한다. \textbf{검증}: 마스킹 패러다임에서 자극 제시 시간을 조작하여, 폐합 판단과 분리 판단의 시간 순서를 측정.

\paragraph{E2. 메타안정성과 지각적 체류 (등급 A).}
T7-Enhanced에 의해, 폐합이 강화된 형성은 더 깊은 에너지 분지에 있으므로, 양안 경쟁(binocular rivalry)이나 다의적 형태(ambiguous figures)에서 더 긴 지각적 체류 시간을 보여야 한다. \textbf{검증}: 폐합 강도(불완전한 Kanizsa 형태의 완성도)를 조작하고 체류 시간을 측정.

\paragraph{E3. $\beta_{\mathrm{crit}}$ 대응물 (등급 B).}
SCC의 위상 전이 임계값은 지각적 조직화의 급격한 전환에 대응할 수 있다. 자극의 대비나 공간 주파수를 연속적으로 변화시킬 때, 형태 지각이 급격하게 출현하는 임계값이 존재해야 한다. \textbf{검증}: 무작위점(random dot) 디스플레이에서 응집도를 연속적으로 변화시키며 형태 검출 임계값 측정.

\subsection{신경영상 예측}

\paragraph{E4. 자기-참조적 피드백 루프 (등급 B).}
SCC의 에너지 최소화는 피드포워드/피드백 처리의 반복을 필요로 한다. EEG/MEG에서 형태 지각 관련 감마 대역 활동은 $\mathrm{Cl}_t$의 반복 적용에 대응하는 재진입(reentrant) 패턴을 보여야 한다. \textbf{검증}: 시간-주파수 분석에서 형태 지각 전후의 재진입 감마 활동 측정.

\paragraph{E5. 진단 벡터의 신경 상관물 (등급 B).}
$\mathbf{d}$의 네 성분(Bind, Sep, Inside, Persist)은 서로 다른 신경 과정에 대응할 수 있다. fMRI에서 Bind는 영역 내 동기화에, Sep는 영역 간 비동기화에, Inside는 위상적 연결 패턴에 대응해야 한다. \textbf{검증}: 다변량 패턴 분석(MVPA)으로 진단 벡터의 네 차원을 신경 활동에서 디코딩.

\paragraph{E6. 형성 경계와 신경 활동 (등급 A).}
SCC의 경계 대역 $\mathrm{Bd}_t$에 해당하는 자극 영역에서 가장 높은 신경 반응이 관찰되어야 한다 (경계에서 에너지 기울기가 최대). \textbf{검증}: 망막 위상 지도(retinotopic mapping)에서 형성 경계에 대응하는 V1/V2 활동 측정.

\subsection{발달심리학 예측}

\paragraph{E7. 발달적 매개변수 변화 (등급 C).}
영아의 지각적 조직화 발달은 SCC 매개변수의 변화로 모델링될 수 있다. 초기 발달에서 $\beta_{\mathrm{bd}}$ (경계 선명도)가 증가하고, 이후 $\lambda_{\mathrm{cl}}$ (폐합 강도)이 증가하는 순서가 예측된다. \textbf{검증}: 선호 주시(preferential looking) 패러다임에서 연령에 따른 형태 분절/완성 능력의 발달 순서 측정.

\paragraph{E8. 대상 영속성과 Persist (등급 B).}
피아제의 대상 영속성은 SCC의 $\mathsf{Persist}$ 진단에 대응한다. 가림(occlusion) 상황에서 형성의 $\mathsf{Persist}$가 문턱값 이상이면 대상이 ``존재한다''고 인식하고, 이하이면 ``사라졌다''고 인식해야 한다. \textbf{검증}: 영아의 가림 실험에서 가림 시간/정도를 조작하고, SCC 모델의 $\mathsf{Persist}$ 예측과 영아의 놀라움 반응을 비교.

\subsection{일반 이론 예측}

\paragraph{E9. 다중 형성의 경쟁 (등급 B).}
$K > 1$ 형성이 $K = 1$ 전역 최소보다 에너지적으로 불리함에도 지각적으로 안정적인 것은 메타안정성 때문이다. 지각적 조직화에서 ``여러 물건'' 인식이 ``하나의 덩어리'' 인식보다 에너지적으로 불리하지만 지각적으로 더 안정적인 상황이 존재해야 한다. \textbf{검증}: 다의적 자극에서 ``여러 물건'' 해석과 ``하나의 덩어리'' 해석의 전환 동역학 측정.

\paragraph{E10. Sep-before-Inside 시간 순서 (등급 A).}
P5의 직접적 인지 대응. 자극이 점진적으로 드러날 때, 내부/외부 판별(Sep)이 형태 분절(Inside)보다 시간적으로 선행해야 한다. \textbf{검증}: 점진적 자극 출현(gradual onset) 패러다임에서 두 판단의 반응 시간 비교.


%% ============================================================
\section{응용: 컴퓨터 비전}
\label{sec:cv-applications}

\subsection{부드러운 분할}

기존 영상 분할은 이진(binary) 출력을 강제한다. SCC의 응집 장 $u : X \to [0,1]$은 본래 \textbf{부드러운 분할}(soft segmentation)을 제공하며, 경계의 모호성, 겹침, 점진적 전이를 자연스럽게 포착한다.

$K$-장 아키텍처를 영상 분할에 적용하면, 각 ``물건''이 하나의 응집 장으로 표현되며, 형성 간 반발이 공간적 분리를 유도하고, 심플렉스 제약이 픽셀당 총 소속을 1 이하로 제한한다.

\subsection{자기-참조 그래프 신경망}

SCC의 핵심 구조---자기-참조적 에너지 최소화---는 그래프 신경망(GNN)의 새로운 아키텍처를 시사한다:
\begin{enumerate}
    \item \textbf{표준 GNN}: 메시지 전달 $\to$ 집약 $\to$ 갱신 (피드포워드)
    \item \textbf{SCC-GNN}: 메시지 전달 $\to$ 자기-참조 에너지 계산 $\to$ 경사 흐름 수렴 $\to$ 출력
\end{enumerate}
SCC-GNN은 반복적 에너지 최소화를 통해 전역적 일관성(global coherence)을 달성하며, 이는 표준 메시지 전달의 수용 영역 제한을 극복한다.


%% ============================================================
\section{응용: 인지 모델링}
\label{sec:cognitive-applications}

\subsection{대상 영속성}

피아제의 대상 영속성은 SCC에서 자연스럽게 모델링된다:
\begin{enumerate}
    \item 가림 전: 형성 $\hat{u}_t$가 존재 ($\mathsf{Bind}, \mathsf{Sep}, \mathsf{Inside}$ 높음)
    \item 가림 중: 직접적 감각 입력 없이, 시간적 전달 $\mathbf{M}_{t \to s}$가 형성을 전파
    \item $\mathsf{Persist}$ 값이 문턱값 이상이면 ``대상이 여전히 존재''로 인식
    \item $\mathsf{Persist}$가 문턱값 이하로 떨어지면 대상 표상 소실
\end{enumerate}

\subsection{전경-배경 분리}

전경-배경(figure-ground) 분리는 SCC의 단일 형성 발견에 정확히 대응한다. 형성의 핵심(Core)이 전경이고, 외부(Ext)가 배경이며, 경계 대역(Bd)이 전이 영역이다. SCC의 독특한 기여는 이 분리가 이진적이 아닌 \emph{연속적}이라는 점이다.

\subsection{주의적 그룹화}

주의(attention)는 SCC에서 체적 제약 $\Sigma_m$의 매개변수 $c$로 모델링할 수 있다:
\begin{itemize}
    \item $c$ 작음: 적은 자원, 하나의 좁은 형성 (집중 주의)
    \item $c$ 큼: 많은 자원, 넓거나 여러 형성 (분산 주의)
    \item $K$-장: 여러 대상에 대한 동시 주의 (분할 주의)
\end{itemize}

% Figure placeholder
\begin{figure}[t]
\centering
\fbox{\parbox{0.85\textwidth}{\centering\vspace{5cm}
\textbf{[그림 14.3]} 응용 영역 개관.\\
좌: 컴퓨터 비전---부드러운 분할, 자기-참조 GNN.\\
중: 인지 모델링---대상 영속성, 전경-배경, 주의.\\
우: 신경과학---진단 벡터의 신경 상관물.
\vspace{0.5cm}}}
\caption{SCC 이론의 세 가지 응용 영역. 각 영역에서 SCC의 수학적 구조가 구체적 문제에 대한 새로운 접근을 제공한다.}
\label{fig:applications}
\end{figure}


%% ============================================================
\section{열린 문제와 미래 방향}
\label{sec:open-problems}

SCC 이론의 현재 형태에서 여러 중요한 문제가 미해결로 남아 있다.

\subsection{연속체 극한}

현재 SCC는 유한 그래프에서 정의된다. $\alpha = \alpha_0 / h^2$ 재조정을 통해 연속 극한의 단서가 있지만(스케일링 각주, T8-Core), 완전한 연속체 SCC---편미분 방정식으로서의 정식화---는 아직 달성되지 않았다.

\begin{itemize}
    \item 폐합 연산자의 연속 극한: 비국소 확산 연산자?
    \item 에너지의 $\Gamma$-수렴: 격자 $\to$ 연속체에서의 수렴 보장은?
    \item 이동 프론트(traveling front) 해의 존재?
\end{itemize}

\subsection{확률적 확장}

현재 SCC는 결정론적이다. 확률적 확장을 통해:
\begin{itemize}
    \item Kramers 장벽 교차에 의한 $K \to K-1$ 전이 속도
    \item 노이즈 유도 상전이와 확률적 공명
    \item 형성의 수명과 안정성의 확률적 특성화
\end{itemize}

\subsection{매개변수 학습}

현재 매개변수는 수동으로 설정되거나 헤시안 정규화로 조정된다. 데이터로부터 매개변수를 학습하는 것은 응용에 필수적이다:
\begin{itemize}
    \item 이중 수준 최적화(bilevel optimization): 상위 수준에서 매개변수, 하위 수준에서 형성
    \item SCC 에너지의 미분 가능성이 역전파를 허용
    \item 주의: 매개변수 학습이 이론의 원초적 구조를 변형시키지 않도록 주의 필요
\end{itemize}

\subsection{신경 구현}

SCC의 에너지 최소화가 실제 신경 회로에서 어떻게 실현될 수 있는가?
\begin{itemize}
    \item 폐합 $\mathrm{Cl}_t$: 재진입(reentrant) 연결의 반복적 활성화?
    \item 에너지 경사 하강: 신경 집단의 율(rate) 동역학?
    \item 위상 전이: 피질 흥분-억제 균형의 분기?
\end{itemize}

이들은 이론적 신경과학과의 깊은 대화를 필요로 하며, SCC의 가장 야심찬 미래 방향이다.

\subsection{연구 로드맵}

열린 문제들을 우선순위로 정리하면:

\begin{table}[h]
\centering
\caption{열린 문제 우선순위 로드맵}
\label{tab:roadmap}
\small
\begin{tabular}{clcc}
\toprule
\textbf{순위} & \textbf{문제} & \textbf{난이도} & \textbf{영향} \\
\midrule
1 & 매개변수 학습 (이중 수준 최적화) & 중 & 매우 높음 \\
2 & 확률적 확장 (Kramers) & 중-고 & 높음 \\
3 & 연속체 극한 (PDE 정식화) & 고 & 매우 높음 \\
4 & 신경 구현 모델 & 고 & 높음 \\
5 & 계층적 SCC & 매우 고 & 혁신적 \\
\bottomrule
\end{tabular}
\end{table}


%% ============================================================
\section{응용: 신호 처리와 신경과학}
\label{sec:signal-neuro-applications}

SCC의 수학적 구조는 시각적 장면 외에도 다양한 신호 처리 문제에 적용될 수 있다.

\subsection{음성 분절}

음성 신호의 분절(segmentation)은 연속적인 음성 스트림에서 단어나 음절의 경계를 찾는 문제이다. SCC 프레임워크로 이를 접근할 수 있다:
\begin{itemize}
    \item \textbf{관계적 지지 공간}: 시간-주파수 표현(스펙트로그램)의 시간-주파수 격자를 그래프 $X$로 구성
    \item \textbf{인접 구조}: 시간적/주파수적 근접성에 기반한 인접 행렬
    \item \textbf{응집 장}: 각 시간-주파수 빈의 ``하나의 음성 단위에 속하는 정도''를 나타내는 연속 장
    \item \textbf{시간적 전달}: 음절 간 전이에서의 구조적 계승
\end{itemize}

기존 음성 분절과의 핵심 차이: SCC는 경계를 미리 정하지 않고, 에너지 최소화를 통해 부드러운 분절을 자동으로 발견한다. 전이 대역의 폭이 음성의 공동 조음(coarticulation)을 자연스럽게 포착한다.

\subsection{EEG/fMRI 응집 분석}

신경영상 데이터에서 ``응집적 뇌 영역''을 발견하는 문제는 SCC의 자연스러운 응용이다:
\begin{itemize}
    \item \textbf{EEG}: 전극 간 동기화(coherence)를 인접 행렬로, 각 전극의 활성화를 초기 장으로 설정. SCC 에너지 최소화가 응집적 전극 클러스터를 발견.
    \item \textbf{fMRI}: 복셀(voxel) 간 기능적 연결을 인접 행렬로 구성. 부드러운 분할이 이진 클러스터링보다 풍부한 네트워크 구조를 포착.
    \item \textbf{시간 역학}: 시간적 전달 $\mathbf{M}_{t \to s}$가 동적 기능 연결(dynamic functional connectivity)의 시간적 진화를 모델링.
\end{itemize}

\paragraph{ICA/NMF와의 비교.} 독립 성분 분석(ICA)이나 비음수 행렬 분해(NMF)와 달리, SCC는:
\begin{enumerate}
    \item 선형 혼합 가정이 불필요
    \item 공간적 구조(인접성)를 명시적으로 활용
    \item 시간적 계승을 내장 (별도의 시간 모델 불필요)
\end{enumerate}

\subsection{신경 집단 분석}

칼슘 영상이나 다전극 어레이(MEA) 데이터에서 신경 집단(neuronal ensemble)을 발견하는 문제:
\begin{itemize}
    \item 뉴런 간 상관 구조를 인접 행렬로 구성
    \item SCC의 $K$-장 아키텍처가 여러 동시 앙상블을 발견
    \item $\mathsf{Persist}$ 진단이 앙상블의 시간적 안정성을 정량화
\end{itemize}

% Figure placeholder
\begin{figure}[t]
\centering
\fbox{\parbox{0.85\textwidth}{\centering\vspace{5cm}
\textbf{[그림 14.4]} 신호 처리와 신경과학 응용.\\
좌: 스펙트로그램 위의 음성 분절---응집 장이 음절 경계를 부드럽게 포착.\\
중: fMRI 기능 연결 그래프 위의 SCC 형성---이진 클러스터링보다 풍부한 구조.\\
우: 신경 앙상블 발견---다전극 데이터에서 $K$-형성으로 동시 활성 뉴런 그룹 식별.
\vspace{0.5cm}}}
\caption{SCC의 신호 처리 및 신경과학 응용. 관계적 지지 공간의 일반성이 시각적 장면 외의 다양한 도메인으로의 확장을 가능하게 한다.}
\label{fig:signal-neuro}
\end{figure}


%% ============================================================
\section{연습 문제}
\label{sec:ch14-exercises}

\begin{enumerate}
    \item \textbf{게슈탈트 매핑 확장.} Pr\"agnanz (간결성) 외에 ``전경-배경 분리'' 법칙을 SCC에 대응시켜라. SCC의 어떤 구성 요소가 전경과 배경의 비대칭성을 포착하는가?

    \item \textbf{프레임워크 비교.} SCC와 예측적 처리(PP)에서 ``놀라움''(surprise)의 대응물은 무엇인가? SCC 에너지의 어떤 항이 PP의 예측 오류에 가장 가까운가? 논거를 제시하라.

    \item \textbf{예측 설계.} E1 (폐합 선행 효과)을 검증하는 심리물리학 실험을 상세히 설계하라. 자극, 과제, 종속 변수, 통제 조건을 명시하라.

    \item \textbf{음성 분절 모델.} 2초 분량의 단순 음성 스펙트로그램(시간 50 빈 $\times$ 주파수 20 빈 = 1000 노드)을 그래프로 구성하라. 시간적/주파수적 인접 구조를 정의하고, SCC 형성을 찾아라. 형성이 음절 경계와 대응하는지 논의하라.

    \item \textbf{SCC-GNN 아키텍처.} 표준 GCN (Graph Convolutional Network)에 SCC의 자기-참조 에너지 최소화를 내장하는 아키텍처를 설계하라. 순방향 패스, 손실 함수, 역전파의 수학적 형태를 기술하라.

    \item \textbf{열린 문제 탐색.} SCC의 연속체 극한에서 폐합 연산자 $\mathrm{Cl}_t$의 연속 유사체를 제안하라. 비국소 확산 방정식의 형태를 쓰고, A1'--A4 공리와의 대응을 논의하라.
\end{enumerate}


%% ============================================================
\section{요약}
\label{sec:ch14-summary}

SCC 이론은 게슈탈트 심리학의 형식화, 예측적 처리/IIT/베이지안 뇌와의 구조적 비교, 10개의 검증 가능한 경험적 예측, 그리고 컴퓨터 비전/인지 모델링에의 응용 가능성을 통해 순수 수학을 넘어선 가치를 입증한다.

특히, 게슈탈트 법칙에 대한 정량적 실현과 Sep-before-Inside 순서의 예측은 SCC의 고유한 기여이다. 열린 문제---연속체 극한, 확률적 확장, 매개변수 학습, 신경 구현---는 이론의 미래 발전 방향을 제시한다.

마지막 장에서는 48개의 Category A 정리를 체계적으로 정리하여, 이론의 수학적 완결성을 종합한다.
